\documentclass{article}
\usepackage{amsmath}
\usepackage{geometry}
\usepackage{todonotes}
\usepackage{enumitem}
\usepackage{booktabs}
\usepackage{graphicx}
\usepackage{float}
\usepackage[utf8]{inputenc}   % jeśli plik jest w~UTF‑8
\usepackage[T1]{fontenc}      % albo [OT4] dla „typowo polskiego” kodowania
\usepackage[polish]{babel}
\usepackage{caption}
\usepackage{pgffor}
\usepackage{listings}
\usepackage{verbatim}
\usepackage{fancyvrb}
\usepackage{tabularx}
\usepackage{ragged2e}

% BIBLIOGRAFIA - POPRAWIONA:
\usepackage[numbers,sort&compress]{natbib}
\bibliographystyle{plainnat}

\geometry{a4paper, margin=1in} %rozmiar papieru

\title{Ocena podobieństwa sygnałów akustycznych z~wykorzystaniem wektorowej reprezentacji cech dźwięku
\\-- od cech akustycznych do trudności percepcyjnej: jak porównać dwa dźwięki?
}


% Optimization of Stimulus Sets in Auditory Perception Tests via Vector-Based Signal Similarity Analysis
\usepackage{authblk}
\author[1,*]{Bartosz Lewandowski, \small{(ORCID:0009-0004-1158-0693)}}
\author[2]{Maciej J. Wróbel, \small{(ORCID:0000-0003-1865-2122)}}

\affil[1]{Department of Clinical Engineering, Academy of Silesia in Katowice}
\affil[2]{Department of Otolaryngology and Oncological Laryngology Collegium Medicum NCU}

% Autor korespondencyjny - definicja przypisu
\affil[*]{Corresponding author: Bartosz Lewandowski, Ph.D., bartosz.lewandowski@akademiaslaska.pl}

\date{}  % Empty date removes it entirely

\begin{document}

\maketitle

\begin{abstract}

    Praca przedstawia opracowanie i~walidację obiektywnej, numerycznej miary podobieństwa bodźców akustycznych. Głównym celem badań było stworzenie metryki opartej na wektorowej reprezentacji cech dźwięku, która pozwoli na ilościową ocenę ryzyka błędnej identyfikacji, złożonych z~pojedyńczych dzwięków próbek wynikającej z~ich bliskości akustycznej, a~nie z~deficytów poznawczych słuchaczy.

    Materiał badawczy stanowiło 24 nagrań odgłosów zwierząt, które opisano za pomocą 16 wymiarowych wektorów cech, obejmujących centroid widmowy, częstotliwość podstawową, współczynnik Peak-to-Average Ratio oraz 13 współczynników cepstralnych w~skali melowej. W~celu ujednolicenia wpływu poszczególnych parametrów zastosowano normalizację metodą Min-Max, a~relacje między próbkami poddano analizie z~wykorzystaniem odległości euklidesowej oraz metody głównych składowych.


    Zidentyfikowano kluczowe parametry różnicujące barwę dźwięku oraz wykazano brak silnej korelacji między centroidem widmowym a~częstotliwością podstawową ($r = 0{,}30$), co potwierdza unikalność informacji niesionej przez poszczególne deskryptory. Walidacja krzyżowa z~wynikami badań klinicznych wykazała zgodność opracowanej metryki z~empiryczną trudnością percepcyjną: para bodźców o~najniższym dystansie euklidesowym ($d = 0{,}714$) była identyfikowana poprawnie jedynie w~$20{,}5\%$ przypadków, podczas gdy para o~wyraźnie większym dystansie ($d = 0{,}942$) osiągnęła  $92{,}9\%$ poprawnych odpowiedzi. Opracowana metodologia stanowi fundament do tworzenia powtarzalnych zestawów bodźców o~kontrolowanym stopniu trudności w~testach percepcji słuchowej.

\end{abstract}


\vspace{1em}

\noindent\textbf{Keywords:}  sound perception; sound analysis, vector analysis, acoustic tests, birds sounds, animal sounds, sound properties;

\section{Wstęp}

Konieczność prowadzenia analiz porównawczych sygnałów akustycznych, w~tym złożonych próbek dźwiękowych, wynika z~różnorodnych potrzeb natury poznawczej, diagnostycznej i~aplikacyjnej. Jednym z~istotnych obszarów zastosowań są badania audiologiczne, w~których porównywanie bodźców dźwiękowych o~kontrolowanych parametrach fizycznych umożliwia obiektywną ocenę funkcjonowania narządu słuchu oraz kondycji zdrowotnej pacjenta. 
Podobne w~swojej naturze bodźce dźwiękowe, niosące różne informacje, wybrzmiewające osobno pozostają „czytelne” w odbiorze, podczas gdy oddziałujące równocześnie (wybrzmiewające jednoczasowo), mogą być niezrozumiałe, ze względu na subtelne podobieństwa lub niewielkie różnice w fizycznej naturze tych bodźców. 

Porównawcza analiza sygnałów akustycznych pozwala również na pogłębianie wiedzy dotyczącej mechanizmów percepcji słuchowej, w~tym relacji pomiędzy fizycznymi właściwościami bodźca a~jego subiektywnym odbiorem. W~tym kontekście szczególnego znaczenia nabierają metody umożliwiające ilościowy opis podobieństwa dźwięków, które stanowią podstawę do badań nad rozróżnialnością bodźców, adaptacją sensoryczną oraz przetwarzaniem informacji akustycznej w~ośrodkowym układzie nerwowym. 
Rozwój narzędzi analitycznych pozwalających na obiektywną i~powtarzalną ocenę podobieństwa sygnałów akustycznych stanowi zatem istotny element interdyscyplinarnych badań łączących akustykę, audiologię, neurofizjologię oraz nauki kognitywne. 

Jednym z~kluczowych wyzwań w~badaniach opartych na analizie bodźców akustycznych jest duża liczba zmiennych wpływających na uzyskiwane wyniki, obejmujących zarówno parametry fizyczne sygnału, jak i~czynniki związane z~warunkami jego rejestracji oraz indywidualnymi cechami percepcyjnymi odbiorcy\cite{Letowski1989SoundQuality}. Zmienność tych parametrów może stanowić istotne ograniczenie metodologiczne, prowadząc do trudności w~jednoznacznej interpretacji wyników oraz obniżenia ich powtarzalności i~wiarygodności. W~kontekście testów audiologicznych oraz badań nad percepcją słuchową szczególnie problematyczne jest rozdzielenie wpływu właściwości samego bodźca od efektów wynikających z~jego prezentacji lub subiektywnej interpretacji. 

Zastosowanie ustrukturyzowanych metod analizy porównawczej, opartych na mierzalnych i~jednoznacznie definiowalnych parametrach akustycznych, umożliwia systematyczne uwzględnienie tej zmienności i~ograniczenie jej wpływu na końcowe wnioski badawcze. Reprezentacja próbek dźwiękowych w~postaci wielowymiarowych wektorów cech pozwala na równoczesną analizę wielu aspektów sygnału, takich jak struktura widmowa, dynamika amplitudowa czy cechy cepstralne, bez konieczności każdorazowego odwoływania się do złożonej i~czasochłonnej oceny fizycznej poszczególnych parametrów. Takie podejście sprzyja standaryzacji procedur analitycznych, zwiększa obiektywność porównań oraz umożliwia szybką i~powtarzalną ocenę podobieństwa sygnałów akustycznych. W~rezultacie analiza wektorowa stanowi efektywne narzędzie wspomagające interpretację wyników badań, pozwalając na bardziej przejrzyste odniesienie obserwowanych różnic i~podobieństw do mierzalnych właściwości bodźców dźwiękowych.

Istnieje szereg metod analizy próbek dźwiękowych, pozwalających na szczegółową charakterystykę fizyczną, opartych na komplementarnym zastosowaniu analiz w~domenie czasu, częstotliwości oraz czasowo-częstotliwościowej~\cite{Letowski1989SoundQuality,Peeters2004CUIDADO}. I~tak, analiza czasowa może obejmować ocenę przebiegu sygnału w~funkcji czasu poprzez wyznaczenie parametrów amplitudowych (wartość skuteczna RMS, dynamika), \textit{Zero-Crossing Rate} (ZCR) do wstępnej klasyfikacji okresowości/szumów~\cite{Gouyon2000}, oraz charakterystykę obwiedni (ADSR: attack, decay, sustain, release) określającą czasowe właściwości ataku i~wybrzmienia~\cite{Peeters2004CUIDADO}. Z~kolei analiza częstotliwościowa wykorzystuje transformatę Fouriera do określenia widma amplitudowego i~fazowego, identyfikacji składowych harmonicznych, \textit{Spectral Centroid} (jasność barwy), \textit{Spectral Flatness} (tonalność) oraz \textit{Spectral Roll-off} (koncentracja energii widma)~\cite{Peeters2004CUIDADO,saitis2020brightness}. 
Natomiast dla sygnałów niestacjonarnych stosuje się Krótkoczasową Transformatę Fouriera (STFT) i~spektrogramy do śledzenia zmian widma w~czasie, uzupełnione deskryptorami psychoakustycznymi jak MFCC i~PLP, modelującymi nieliniowość ludzkiego słuchu poprzez banki filtrów melowych i~krytycznych~\cite{Hermansky1990,tracey2023interpretable_mfcc}.
Tak kompleksowe podejście analityczne umożliwia obiektywną ocenę fizycznych i~percepcyjnych właściwości określonych próbek dźwiękowych, co może stanowić podstawę dalszych analiz podobieństwa ocenianych próbek~\cite{Peeters2004CUIDADO}.

Porównanie próbek dźwiękowych w~oparciu o~ich fizyczne składowe, tj. właściwości fizyczne bodźca dźwiękowego, oraz próby wskazania czy określone składowe różnych bodźców dźwiękowych odpowiadają za trudności w~rozumieniu przekazywanej przez każdą z~tych próbek informacji, ze względu na złożoność procesów analitycznych jest dużym wyzwaniem. Stąd istnieje potrzeba opracowania prostszych, parametrycznych, bardziej czytelnych algorytmów analizy poziomu podobieństw lub/i~różnic między bodźcami dźwiękowymi poddanymi ocenie.

Niniejsza praca podejmuje temat selekcji cech fizycznych bodźców akustycznych, dalej przyporządkowania i~obliczenia porównywalnych wielkości wektorowy dla poszczególnych cech fizycznych próbek dźwiękowych, w~celu weryfikacji hipotezy badawczej zakładającej, że odległość euklidesowa między znormalizowanymi wektorami określonych cech akustycznych próbek dźwiękowych, stanowi obiektywną miarę podobieństwa (lub różnicy) porównywanych bodźców dźwiękowych~\cite{MadhavanWrobel2024MLT}, która koreluje z~wynikami ich rozróżnienia w~testach klinicznych. W~szczególe, autorzy podjęli próbę weryfikacji założenia: że przyczyną błędnej identyfikacji podobieństwa pary bodźców dźwiękowych odtwarzanych podczas badania klinicznego może być „mały” dystans w~przestrzeni euklidesowej wcześniej wybranych i~zdefiniowanych cech [wybranych i~zdefiniowanych dla potrzeb tej analizy]  w~przeciwieństwie do prawidłowej identyfikacji podobieństwa bodźców dźwiękowych dla pary o~„dużym” dystansie w~przestrzeni euklidesowej. 


\section{Materiały i~metodyka}
\label{section:mat_met}

Głównym założeniem analizy było wypracowanie konkretnej miary dla złożonych bodźców dźwiękowych, która pozwoli matematycznie ocenić, jak bardzo poszczególne sygnały - bodźce dźwiękowa są do siebie podobne. Innymi słowy, celem pracy było Część A. opracowanie metryki podobieństwa akustycznego, bazującego na odległościach między znormalizowanymi wektorami cech akustycznych próbek dźwiękowych, a~następnie Część B. zweryfikowanie metryk dźwiękowych z~wynikami doświadczeń klinicznych percepcji złożonych bodźców dźwiękowych.

Materiał badawczy stanowiły próbki dźwiękowe reprezentujące odgłosy popularnych zwierząt. Nagrania przygotowane w~formie plików MP3 (Moving Picture Experts Group Layer-3), próbkowanych z~częstotliwością 44 kHz. Analizowany zbiór 24 nagrań cechował się zróżnicowaną długością trwania -- od 3,2~s (\textit{hawk.mp3}) do 162,3~s (\textit{duck.mp3}), przy średniej 32,9~s i~medianie 9,5~s, w~trybie nagrania mono; łączny rozmiar zbioru wyniósł około 13,7~MB. 

Łącznie przygotowano 24 pliki dźwiękowe reprezentujące różne gatunki zwierząt: z~czego 10 plików było już wykorzystanych w~pracy Madhavan et al\cite{MadhavanWrobel2024MLT}:chicken.mp3, cow.mp3, dog.mp3, donkey.mp3, elephant.mp3, frog.mp3, horse.mp3, lion.mp3, monkey.mp3 sheep.mp3 oraz dodatkowe 14 dźwięków zwierząt wybranych losowo z biblioteki BBC Sound Effects\cite{bbcrewind_soundeffects}:  blackbird.mp3, chaffinch.mp3, cow.mp3, donkey.mp3, elephant.mp3, frog.mp3, great\_tit.mp3, horse.mp3, monkey.mp3, sheep.mp3, tundra\_swan.mp3, wolf1.mp3, cat.mp3, chicken.mp3, dog.mp3, duck.mp3, frog2.mp3, gorilla.mp3, hawk.mp3, lion.mp3, owl.mp3, tawny\_owl.mp3, wild\_boar.mp3, wolf2.mp3.

Zestaw cech dobrano zgodnie z~zasadą możliwie oszczędnego, lecz komplementarnego opisu sygnału, tak aby pokryć trzy względnie niezależne domeny percepcyjno-akustyczne przy zachowaniu interpretowalności każdego wymiaru: barwę opisano centroidem widmowym (najlepiej ugruntowanym korelatem percepcyjnej „jasności” brzmienia~\cite{Grey1977Timbre,Lakatos2000Timbre,saitis2020brightness}) oraz zestawem 13 współczynników MFCC stanowiących standardowy, zwarty opis kształtu obwiedni widmowej~\cite{Peeters2011TimbreToolbox,Siedenburg2017TimbreSimilarity,DeCamargo2017ProtaxSound}, przyjętych w~liczbie zgodnej z~powszechną konwencją i~domyślną konfiguracją biblioteki \textit{librosa}~\cite{mcfee2015librosa}; wysokość dźwięku reprezentowano częstotliwością podstawową (F0), której niezależność od barwy potwierdza słaba korelacja z~centroidem ($r = 0{,}30$); dynamikę czasową ujęto parametrem PAR, pozwalającym odróżnić sygnały o~zbliżonej wysokości i~barwie, lecz odmiennym charakterze impulsowym. Za takim doborem przemawiały ponadto względy praktyczne: jednoznaczna interpretacja każdej cechy (istotna przy zastosowaniu metryki do audytu materiału bodźcowego), niewielki wymiar wektora adekwatny do liczebności zbioru (24~nagrania), dostępność wszystkich deskryptorów w~standardowych, otwartych narzędziach ekstrakcji oraz zgodność z~reprezentacją stosowaną w~pracy Madhavan i~wsp.~\cite{MadhavanWrobel2024MLT}, którą niniejsze badanie rozszerza. Przyjęto zatem celowo prostą i~interpretowalną reprezentację bazową, stanowiącą punkt odniesienia dla bardziej rozbudowanych zestawów cech (por. Dyskusja).

Każdy dźwięk opisano za pomocą zestawu 16 cech akustycznych obejmujących: A. centroid widmowy opisujący „jasność” dźwięku, B. częstotliwość podstawową (\textit{F0}), C. Peak to Average Ratio (\textit{PAR}); D. Mel Frequency Cepstral Coefficients (\textit{MFCC}) -13 cech opisujących dynamikę oraz barwę, wyekstrahowanych przy użyciu biblioteki \textit{librosa}~\cite{mcfee2015librosa} (Tab.~\ref{tab:parametry-akustyczne}).
\newcolumntype{L}{>{\RaggedRight\arraybackslash}X}
\begin{table}[H]
\centering
\scriptsize
\renewcommand{\arraystretch}{1.0}
\setlength{\aboverulesep}{1pt}
\setlength{\belowrulesep}{1pt}
\begin{tabularx}{\textwidth}{@{}>{\bfseries}p{3.2cm}LL@{}}
\toprule
\textbf{Parametr} & \textbf{Opis techniczny} & \textbf{Interpretacja akustyczna} \\
\midrule
Centroid widmowy
& Parametr opisujący środek ciężkości widma sygnału dźwiękowego, określający częstotliwość, wokół której skupiona jest energia sygnału.
& Wyższe wartości oznaczają dominację wysokich częstotliwości (dźwięk jasny, ostry, klarowny). Niższe wartości wskazują na przewagę niskich częstotliwości (dźwięk stłumiony, ciemny). \\
\midrule
Częstotliwość podstawowa (F0; Hz)
& Najniższa składowa częstotliwościowa złożonej próbki dźwiękowej, odpowiadająca wysokości dźwięku.
& Niższe F0 oznacza dźwięk bardziej basowy, wyższe F0 -- wyższy, bardziej ,,piszczący''. \\
\midrule
Peak-to-Average Ratio (PAR; dB)
& Stosunek wartości szczytowej sygnału do jego wartości średniokwadratowej (RMS), wyrażony w~decybelach.
& Niskie wartości wskazują na stabilną dynamikę (np. mruczenie, przedłużone dźwięki tonalne). Wysokie wartości świadczą o~obecności krótkotrwałych impulsów i~transjentów. Umożliwia rozróżnianie sygnałów o~podobnym F0, lecz odmiennej dynamice. \\
\midrule
MFCC 1 (Energia)
& Pierwszy współczynnik cepstralny opisujący całkowitą energię sygnału.
& Wysokie wartości odpowiadają głośnym dźwiękom, niskie -- cichym sygnałom, szeptowi lub ciszy. \\
\midrule
MFCC 2 (Jasność/Ciemność)
& Współczynnik opisujący rozkład energii między niskimi i~wysokimi częstotliwościami.
& Określa, czy dźwięk jest bardziej jasny (wysokotonowy) czy ciemny (niskotonowy). \\
\midrule
MFCC 3--4 (Kształt widma)
& Współczynniki opisujące strukturę środkowej części widma oraz obecność charakterystycznych rezonansów.
& Pozwalają odróżniać dźwięki o~tej samej wysokości, ale różnej barwie, np. samogłoski o~podobnym F0. \\
\midrule
MFCC 5--8 (Struktura formantowa i~rezonanse)
& Parametry związane z~formantami i~charakterystycznymi rezonansami źródła dźwięku.
& Odpowiadają za indywidualny charakter brzmienia, umożliwiają identyfikację instrumentu lub cech aparatu głosowego. Są kluczowe w~systemach rozpoznawania mowy. \\
\midrule
MFCC 9--13 (Mikrostruktura sygnału)
& Współczynniki opisujące wysokoczęstotliwościowe szczegóły i~drobne fluktuacje widma.
& Określają stopień gładkości, szorstkości lub chropowatości dźwięku. Pomagają odróżniać bardzo podobne próbki dźwiękowe i~zwiększają realizm cyfrowego odwzorowania. \\
\bottomrule
\end{tabularx}
\caption{Zestaw cech akustycznych wybranych do charakterystyki złożonego bodźca dźwiękowego.}
\label{tab:parametry-akustyczne}
\end{table}


% Dla każdego pliku dźwiękowego wykonano analizę cech, uzyskując wartość numeryczną dla każdej z~16 analizowanych cech. 

% \begin{enumerate}[label=\Alph*.]
%     \item \textit{Centroid widmowy} - parametr opisujący środek ciężkości widma sygnału dźwiękowego, określający częstotliwość, wokół której skupiona jest energia sygnału;  wyższe wartości centroidu wskazują na dominację wysokich częstotliwości i~są percepcyjnie związane z~„jasnym”, ostrym i~klarownym brzmieniem, natomiast niższe wartości odpowiadają przewadze niskich częstotliwości i~są odbierane jako dźwięk stłumiony \cite{saitis2020brightness}.
%     \item \textit{Częstotliwość podstawowa} (\textit{F0}; Hz) - najniższa składowa częstotliwościowa złożonej próbki dźwiękowej, potocznie utożsamiana z~"wysokością" dźwięku,  tym, co odróżnia głęboki bas od przenikliwego piszczałki; im niższe F0, tym dźwięk jest niższy i~bardziej basowy (np. ryk lwa), im wyższe  tym wyższy i~bardziej cienki (np. świergot ptaka); F0 jest bezpośrednio powiązana z~tempem drgań struny głosowej lub innego źródła dźwięku i~stanowi jeden z~głównych wyznaczników tożsamości brzmienia danego gatunku, co potwierdzają badania nad zależnością F0 od budowy anatomicznej oraz stanu emocjonalnego zwierząt~\cite{f0_garcia2017acoustic_allometry,f0_briefer2012vocal_expression_emotions}.
%     \item \textit{Peak-to-Average Ratio }(\textit{PAR; dB}) - parametr określający relację pomiędzy maksymalną amplitudą sygnału a~jego średnią wartością amplitudową lub energetyczną, wyrażony w~decybelach; wskazuje na obecność krótkotrwałych, intensywnych pików wyraźnie przekraczających jego średni poziom; niskie wartości PAR są charakterystyczne dla dźwięków stabilnych, takich jak mruczenie lub odgłosy kontaktowe; PAR umożliwia rozróżnianie sygnałów o~tej samej częstotliwości podstawowej (F0), lecz odmiennej dynamice i~charakterze akustycznym.
%     \item \textit{Mel Frequency Cepstral Coefficients} (\textit{MFCC}) - na potrzeby analizy zestaw 13 parametrów opisujących sygnał dźwiękowy w~domenie cepstralnej, stanowiących numeryczną reprezentację jego cech akustycznych; współczynniki MFCC wyznaczane są na podstawie skali melowej (nieliniowej skali częstotliwości modelującej sposób, w~jaki ludzkie ucho odbiera dźwięk: ucho jest bardziej czułe na różnice między dźwiękami niskimi niż wysokimi, dlatego skala melowa "zagęszcza" niskie częstotliwości i~"rozciąga" wysokie, oddając rzeczywistą percepcję słuchową zamiast fizycznego pomiaru w~hercach) i~odzwierciedlają rozkład energii sygnału w~pasmach istotnych fizjologicznie, w~przeciwieństwie do surowego zapisu fali akustycznej; proces ekstrakcji obejmuje analizę widmową sygnału, filtrację z~użyciem banku filtrów melowych oraz transformację kosinusową (DCT), a~jego wynikiem jest wektor kilkunastu współczynników opisujących barwę dźwięku  \cite{tracey2023interpretable_mfcc, lou2021feature}:
%     \begin{enumerate}[label=\roman*.]
%     % Oryginalna treść (ton potoczny) przed ujednoliceniem rejestru stylistycznego:
%     % \item MFCC 1 (Energia): To informacja o tym, jak głośny jest dźwięk. Jeśli ta liczba jest wysoka, mamy do czynienia z potężnym sygnałem; jeśli niska - z ciszą lub szeptem.
%     \item MFCC 1 (Energia): współczynnik odzwierciedlający całkowitą energię sygnału. Wysokie wartości wskazują na sygnał o~dużej głośności, natomiast wartości niskie odpowiadają dźwiękom cichym, zbliżonym do szeptu lub ciszy.

%     % Oryginalna treść (ton potoczny) przed ujednoliceniem rejestru stylistycznego:
%     % \item MFCC 2 (Jasność/Ciemność): To najważniejszy parametr barwy. Mówi nam, czy w dźwięku przeważają niskie tony (np. mruk basu), czy wysokie (np. pisk fletu). To fundament balansu tonalnego.
%     \item MFCC 2 (Jasność/Ciemność): kluczowy parametr opisujący barwę dźwięku, wskazujący na przewagę niskich częstotliwości (brzmienie ciemne) lub wysokich (brzmienie jasne). Stanowi podstawowy wyznacznik balansu tonalnego sygnału.

%     % Oryginalna treść (ton potoczny) przed ujednoliceniem rejestru stylistycznego:
%     % \item MFCC 3 i 4 (Kształt widma): Odpowiadają za to, czy dźwięk jest "pusty" w środku, czy może ma dużo "nosowości". To dzięki nim możliwe jest odróżnienie od siebie samogłosek o tej samej wysokości tonu, ponieważ wyłapują one specyficzne wzmocnienia w środkowym paśmie częstotliwości.
%     \item MFCC 3 i~4 (Kształt widma): opisują rozkład energii w~środkowym paśmie częstotliwości, w~tym obecność charakterystycznych wzmocnień widmowych (np. rezonansu nosowego). Umożliwiają rozróżnienie dźwięków o~tej samej wysokości tonu, lecz odmiennej barwie, takich jak różne samogłoski.

%     % Oryginalna treść (ton potoczny) przed ujednoliceniem rejestru stylistycznego:
%     % \item MFCC 5-8: Środkowa grupa współczynników, obejmująca numery od 5 do 8, zajmuje się strukturą formantową i specyficznymi rezonansami. w tej fazie dźwięk przestaje być tylko "jasny" lub "ciemny", a zaczyna nabierać swojego charakteru. Te parametry wychwytują subtelne zmiany w barwie, które pozwalają zidentyfikować materiał, z którego wykonany jest instrument, lub specyficzny układ krtani mówcy. Jest to etap, na którym surowe dane spektralne zamieniane są na informacje o unikalnej tożsamości brzmienia, co ma kluczowe znaczenie w systemach rozpoznawania mowy.
%     \item MFCC 5-8 (Struktura formantowa i~rezonanse): środkowa grupa współczynników opisująca strukturę formantową oraz specyficzne rezonanse źródła dźwięku. Parametry te wychwytują subtelne zmiany barwy pozwalające zidentyfikować materiał, z~którego wykonany jest instrument, lub indywidualny układ aparatu głosowego mówcy, co ma kluczowe znaczenie w~systemach rozpoznawania mowy.

%     % Oryginalna treść (ton potoczny) przed ujednoliceniem rejestru stylistycznego:
%     % \item MFCC 9-13: Zestawienie zamykają współczynniki od 9 do 13, które w literaturze przedmiotu często określa się mianem parametrów wysokiej częstotliwości lub mikrostruktury. Nie niosą one ze sobą informacji o głośności czy dominującym tonie, lecz skupiają się na najdrobniejszych fluktuacjach i teksturze sygnału. Opisują one to, czy dźwięk jest gładki i czysty, czy może szorstki, chropowaty i pełen nieregularności. Choć ich wpływ na ogólny obraz wydaje się mniejszy niż w przypadku pierwszych współczynników, to właśnie ta końcowa piątka decyduje o realizmie cyfrowego odwzorowania i pozwala na odróżnienie bardzo podobnych do siebie próbek dźwiękowych.
%     \item MFCC 9-13 (Mikrostruktura sygnału): współczynniki określane w~literaturze przedmiotu jako parametry wysokiej częstotliwości lub mikrostruktury sygnału. Nie niosą informacji o~głośności ani dominującym tonie, lecz opisują najdrobniejsze fluktuacje widma, decydujące o~gładkości bądź chropowatości brzmienia. Mimo relatywnie niewielkiego wpływu na ogólny obraz sygnału, umożliwiają odróżnienie bardzo podobnych do siebie próbek dźwiękowych oraz zwiększają realizm cyfrowego odwzorowania.
    
%     % \item MFCC 5–13 (detale): To coraz drobniejsze szczegóły barwy. Można je porównać do tekstury materiału,  czy dźwięk jest gładki, czy szorstki i~chropowaty. Im wyższy numer parametru, tym mniejszy i~bardziej subtelny detal opisuje.
%     \end{enumerate}
% \end{enumerate}

Ze względu na znaczące różnice w~jednostkach i~skalach mierzonych parametrów (np. częstotliwość w~Hz osiągająca wartości rzędu kilku tysięcy, wobec współczynników MFCC o~znacznie mniejszej amplitudzie), w~celu eliminacji dominacji parametrów o~wysokich wartościach bezwzględnych (np. jak Centroid) nad pozostałymi cechami podczas obliczeń statystycznych, wszystkie wyekstrahowane cechy zostały poddane transformacji do wspólnego przedziału  przy użyciu klasy \textit{MinMaxScaler} z~biblioteki s\textit{cikit-learn} w~wersji 1.7.2~\cite{pedregosa2011scikit}. 
Obliczony dla każdej próbki dźwiękowej wektor cech akustycznych powstał w wyniku połączenia składowych, uśrednionych wartości dla każdej z~obliczanej cechy, w~sumie 16 cech.


W~celu redukcji wymiarowości wielowymiarowego zbioru cech akustycznych oraz umożliwienia ich czytelnej wizualizacji zastosowano analizę głównych składowych (\textit{Principal Component Analysis, PCA}) ~\cite{salimi2025evaluating, liu2023voiceprint}. Punkt wyjścia stanowiła macierz danych, w~której każda próbka dźwiękowa była opisana wektorem cech obejmującym parametry widmowe i~cepstralne, w~tym współczynniki MFCC oraz wybrane miary energetyczne. 
Przed właściwą analizą dane poddano centrowaniu oraz standaryzacji w~celu wyeliminowania wpływu różnic skali pomiędzy poszczególnymi cechami. Następnie obliczono macierz kowariancji, która posłużyła do wyznaczenia wartości i~wektorów własnych definiujących nowe, ortogonalne osie zmienności danych. Wektory własne uporządkowano malejąco według odpowiadających im wartości własnych, a~do dalszej analizy wybrano dwie pierwsze główne składowe (PC1 i~PC2), wyjaśniające największą część całkowitej wariancji zbioru. Ostatecznie próbki dźwiękowe zostały rzutowane na płaszczyznę wyznaczoną przez PC1 i~PC2, co umożliwiło ich rozmieszczenie na dwuwymiarowym wykresie oraz dalszą analizę podobieństw i~dystansów pomiędzy sygnałami w~zredukowanej przestrzeni cech.

Finalnie, nałożenie na siebie obu tych wymiarów pozwala sprowadzić skomplikowany, 16-wymiarowy model audio do czytelnego punktu na płaszczyźnie, dając wgląd w~obiektywne podobieństwa kliniczne badanych próbek.

Wektory cech plików audio porównano za pomocą miary odległości euklidesowej (najkrótszy, mierzony w~linii prostej dystans między dwoma punktami w przestrzeni cech) która w analizie dźwięku pozwala matematycznie określić stopień podobieństwa między sygnałami.


\section{Wyniki}

\subsection{Część A: Analiza statystyczna wektorów cech akustycznych}

Na podstawie przeprowadzonej analizy, opisanej w~Sekcji~\ref{section:mat_met}, uzyskano wyniki, które wskazały występowanie istotnych różnic w~analizowanych parametrach. Wyniki przed normalizacją i~po normalizacji metodą Min-Max dla wszystkich 24 próbek włączonych do analizy:

\begin{table}[h]
\centering
\begin{tabular}{lrrrrrrrr}
\toprule
 parameter & min & max & avg & median & std dev & avg & median & std dev\\
          & &  &   &  &  &(norm) &(norm) &(norm) \\
\midrule
Centroid & 1330.571 & 5692.735 & 3084.120 & 2971.778 & 1241.343 & 0.402 & 0.376 & 0.285 \\
F0 & 437.437 & 2420.861 & 1073.498 & 973.540 & 526.305 & 0.321 & 0.270 & 0.265 \\
PAR & 10.051 & 27.325 & 18.150 & 17.800 & 3.753 & 0.469 & 0.449 & 0.217 \\
MFCC\_1 & -483.222 & -82.167 & -303.938 & -295.563 & 100.673 & 0.447 & 0.468 & 0.251 \\
MFCC\_2 & 33.491 & 162.440 & 98.238 & 107.385 & 39.073 & 0.502 & 0.573 & 0.303 \\
MFCC\_3 & -128.249 & 41.666 & -38.534 & -37.380 & 38.780 & 0.528 & 0.535 & 0.228 \\
MFCC\_4 & -33.489 & 44.812 & 6.366 & 5.713 & 19.912 & 0.509 & 0.501 & 0.254 \\
MFCC\_5 & -57.688 & 18.136 & -10.908 & -9.960 & 17.370 & 0.617 & 0.629 & 0.229 \\
MFCC\_6 & -23.840 & 29.992 & 2.916 & 5.256 & 13.409 & 0.497 & 0.541 & 0.249 \\
MFCC\_7 & -30.839 & 8.275 & -9.227 & -7.591 & 10.084 & 0.553 & 0.594 & 0.258 \\
MFCC\_8 & -23.244 & 43.015 & 3.838 & 5.817 & 14.525 & 0.409 & 0.439 & 0.219 \\
MFCC\_9 & -23.692 & 18.037 & -3.056 & -2.227 & 11.948 & 0.495 & 0.514 & 0.286 \\
MFCC\_10 & -11.306 & 15.074 & 2.882 & 2.699 & 7.006 & 0.538 & 0.531 & 0.266 \\
MFCC\_11 & -22.783 & 17.142 & -0.043 & -1.399 & 10.130 & 0.570 & 0.536 & 0.254 \\
MFCC\_12 & -10.031 & 13.400 & 1.291 & 1.916 & 5.766 & 0.483 & 0.510 & 0.246 \\
MFCC\_13 & -11.244 & 12.754 & 0.405 & 0.247 & 5.676 & 0.485 & 0.479 & 0.237 \\

\bottomrule
\end{tabular}
\caption{Ogólna charakterystyka zbioru dzwięków dla wszystkich 24 próbek włączonych do analizy}
\end{table}


Wyniki znormalizowane metodą Min-Max. Wynik wskazuje, że średnia wartość oscyluje 
w~okolicach $0,4 - 0,55$; oznacza to, że dane są zbalansowane; nie ma sytuacji, 
w~której większość próbek ma skrajnie niskie (blisko $0$) lub skrajnie wysokie 
(blisko $1$) wartości.


\subsubsection{Centroid widmowy (Spectral Centroid)}

Zakres analizy parametru \textit{Spectral Centroid} kształtował się od wartości $5692,735$~Hz (dla pliku \textit{frog.mp3}) do $1330,571$~Hz (dla pliku \textit{gorilla.mp3}). Wartości średnie i~mediana dla analizowanej grupy wyniosły odpowiednio $3084,120$~Hz oraz $2971,778$~Hz, przy odchyleniu standardowym na poziomie $1241,343$~Hz. Większość próbek koncentruje się w~niższych rejestrach jasności, co potwierdza znormalizowana wartość mediany ($\text{med} = 0,376$). Próbki \textit{frog.mp3} ($c = 5692,735$~Hz) oraz \textit{tawny\_owl.mp3} ($c = 5257,504$~Hz) znacząco podnoszą średnią wartość centroidu i~stanowią wartości odstające względem reszty zbioru.
Próbki te oddalone są od kolejnego najwyższego wyniku w~zbiorze o~ponad $1200$~Hz (w~przypadku \textit{tawny\_owl.mp3}), podczas gdy różnica między samym liderem (\textit{frog.mp3}) a~drugim w~kolejności wynikiem wynosi $435,231$~Hz. Tak gwałtowny wzrost wartości w~górnym kwartylu, w~zestawieniu z~odchyleniem standardowym wynoszącym $1241,343$~Hz, jednoznacznie definiuje te dwa nagrania jako rzadkie przypadki o~skrajnie wysokiej jasności spektralnej. Rozkład wartości centroidu dla wszystkich analizowanych próbek przedstawiono na Rys.~\ref{fig:centroid}.

\begin{figure}[H]
    \centering
    \includegraphics[width=0.7\textwidth]{images/chart-centroid.jpeg}
    \caption{Centroid widomowy dla poszczególnych plików dźwiękowych, słupki zielone z~deseniem oznaczają dzwięki pochodzące z~badania\cite{MadhavanWrobel2024MLT}.}
    \label{fig:centroid}
\end{figure}

\subsubsection{Częstotliwość podstawowa (F0)}

Najwyższe wartości: 2420.861 Hz  i~2240.552 Hz odnotowano w~plikach \textit{hawk.mp3} oraz \textit{blackbird.mp3}, natomiast najniższe wartości 437.437 Hz, 456.731 Hz w~plikach \textit{lion.mp3} oraz \textit{tawny\_owl.mp3}; średnia i~mediana dla całej grupy wynosiła odpowiednio (\textit{mean=1073.498Hz}) oraz (\textit{median=973.540Hz}).\\
Analizując dane znormalizowane, F0 ma średnią (\textit{mean=0.321}), co oznacza, że większość próbek oscyluje w~dolnym zakresie wysokości dźwięku w~porównaniu do najwyższej znalezionej tam wartości. Jednocześnie niskie odchylenie standardowe (\textit{std = 0.265}) wskazuje, że próbki są do siebie dość zbliżone pod względem wysokości.
Średnia (\textit{mean=0.321}) jest wyższa od mediany (\textit{median=0.270}) co oznacza, że w~zbiorze jest kilka wysokich dźwięków, które podnoszą średnią, podczas gdy dla większości próbek obliczona wartość F0 jest niska. Zakres (\textit{0.0 – 1.0}) wskazuje, że w~zbiorze jest pełna rozpiętość danych. Wartości częstotliwości podstawowej dla poszczególnych próbek przedstawiono na Rys.~\ref{fig:freq}.

\vspace{1em}

\begin{figure}[H]
    \centering
    \includegraphics[width=0.7\textwidth]{images/chart-f0.jpeg}
    \caption{Częstotliwość podstawowa dla poszczególnych plików dźwiękowych, słupki zielone z~deseniem oznaczają dzwięki pochodzące z~badania \cite{MadhavanWrobel2024MLT}.}
    \label{fig:freq}
\end{figure}

\subsubsection{PAR (Peak-to-Average Ratio)}
Zakres analizy parametru PAR kształtował się od wartości 27.325 dB (dla pliku o~najwyższej dynamice, czyli \textit{duck.mp3}) do wartości 10.051 dB (\textit{dog.mp3}), wartości średnie i~mediana dla analizowanej grupy wyniosły odpowiednio: (mean=18.150 dB), (\textit{median=17.800dB}) oraz odchylenie standardowe (\textit{std = 3.753}).
Analizując dane znormalizowane, parametr PAR, którego średnia (\textit{avg = 0.469}) plasuje się niemal w~połowie skali, wykazuje największą stabilność spośród badanych cech fizycznych. Świadczy o~tym stosunkowo niskie odchylenie standardowe (\textit{std = 0.217}). Wartość mediany (mediana = 0.449) jest bardzo zbliżona do średniej, co sugeruje równomierny rozkład współczynnika szczytu w~całym zbiorze 24 próbek, bez ekstremalnych anomalii. Choć wybrane pliki dominują w~górnym rejestrze skali, to większość danych skupiona jest wokół punktu środkowego, co oznacza, że nagrania mają zbliżoną strukturę dynamiki i~nasycenia sygnału. Rozkład wartości PAR dla analizowanych plików przedstawiono na Rys.~\ref{fig:par}.


\vspace{1em}

\begin{figure}[H]
    \centering
    \includegraphics[width=0.7\textwidth]{images/chart-par.jpeg}
    \caption{PAR dla poszczególnych plików dźwiękowych, słupki zielone z~deseniem oznaczają dzwięki pochodzące z~oryginalnego badania.}
    \label{fig:par}
\end{figure}

\subsubsection{Współczynniki MFCC}

 Zakres analizy pierwszego współczynnika (MFCC\_1), odpowiedzialnego za ogólną energię widmową, kształtował się od $-82.167$ (\textit{gorilla.mp3}) do $-483.222$ (\textit{frog.mp3}), przy średniej wynoszącej $-303.938$. Tak duże odchylenie standardowe ($100.673$) wskazuje, że próbki różnią się między sobą głównie głośnością i~całkowitą mocą sygnału.

W~obszarze charakterystyki spektralnej dominują współczynniki MFCC\_2 (avg = $98.238$) oraz MFCC\_3 (avg = $-38.534$), które niosą najwięcej informacji o~"cieple" lub "ostrości" dźwięku. Przewaga MFCC\_2 sugeruje silnie zdefiniowane niższe pasma częstotliwości. Z~kolei wartości pozostałych współczynników (MFCC\_4 - 13) oscylują bliżej zera (od avg = $-10.908$ do avg = $6.366$ ), co wskazuje na wygładzenie wyższych składowych barwy, typowe dla naturalnych odgłosów zwierząt.

Normalizacja umożliwiła bezpośrednie porównanie wszystkich 13 cech, wykazując wysoki poziom zbalansowania zbioru. Większość współczynników (10 z~13) plasuje się w~przedziale średnich $0.45$-$0.55$, co oznacza, że zbiór jest statystycznie jednorodny, pozbawiony dominującego parametru, który mógłby maskować pozostałe cechy. Najsilniejszym wyróżnikiem okazał się MFCC\_5 z~najwyższą średnią znormalizowaną (avg = $0.617$), to w~tym paśmie kryją się unikalne cechy barwowe definiujące badany zestaw. Największą stabilność wykazuje natomiast MFCC\_8, którego najniższe odchylenie standardowe (std = $0.219$) czyni go wspólnym mianownikiem dla wszystkich 24 nagrań. Parametr MFCC\_9 wyróżnia się wysoką wartością odchylenia standardowego (std = $0.286$), co wskazuje, że informacja przenoszona przez te wartości ma wpływ na próbki dzwiękowe. Odchylenia standardowe poszczególnych współczynników MFCC przedstawiono na Rys.~\ref{fig:stddev}.

\vspace{1em}

\begin{figure}[H]
    \centering
    \includegraphics[width=0.6\textwidth]{images/mfcc_std_dev_radar.jpeg}
    \caption{Odchylenie standardowe dla parametrów MFCC}
    \label{fig:stddev}
\end{figure}


Analiza porównawcza średniej i~mediany wykazała ich wysoką zbieżność we wszystkich 13 współczynnikach, różnice te są minimalne i~zazwyczaj nie przekraczają wartości $0.05$ (np. dla MFCC\_10 avg = $0.538$ vs mediana = $0.531$). Świadczy to o~symetrycznym rozkładzie barwy w~całym zbiorze danych. Można na tej podstawie wnioskować, że nagrania zostały wykonane w~zbliżonych warunkach akustycznych lub przy użyciu urządzeń o~pokrewnej charakterystyce częstotliwościowej, co eliminuje występowanie ekstremalnych anomalii statystycznych w~obszarze psychoakustyki.

Analiza korelacji została użyta jako test redundancji cech, jeśli dwie zmienne są silnie skorelowane, to w~metryce opartej o~odległość euklidesową mogą w~praktyce, w~próbce dzwiękowej,  dublować informację i~wzmacniać jeden wymiar różnic, kosztem drugiego. Niska korelacja oznacza, że cechy opisują odmienne aspekty sygnału i~ich łączne użycie zwiększa rozróżnialność podobieństwa akustycznego (np. niezależność centroidu od F0 oznacza, że barwa/jasność nie wynika bezpośrednio z~wysokości tonu).

\subsubsection{Analiza korelacji parametrów audio}
Macierz korelacji jest tabelarycznym zestawieniem współczynników korelacji obliczonych dla każdej pary analizowanych parametrów, opisujących siłę i~kierunek liniowej zależności między nimi. Współczynnik przyjmuje wartości z~przedziału od $-1$ do $+1$: wartości bliskie $+1$ oznaczają silną zależność dodatnią (wzrost jednej cechy współwystępuje ze wzrostem drugiej), wartości bliskie $-1$ -- silną zależność ujemną (wzrost jednej cechy współwystępuje ze spadkiem drugiej), natomiast wartości zbliżone do zera wskazują na brak liniowego związku, czyli wzajemną niezależność cech. W~kontekście niniejszego badania analiza macierzy korelacji pełni funkcję testu redundancji: cechy silnie skorelowane niosą w~znacznym stopniu tę samą informację i~w~metryce opartej na odległości euklidesowej mogą sztucznie wzmacniać jeden wymiar różnic kosztem pozostałych, podczas gdy cechy słabo skorelowane opisują odmienne aspekty sygnału i~ich łączne uwzględnienie zwiększa zdolność różnicowania podobieństwa akustycznego. Identyfikacja tych zależności pozwala zatem ocenić, czy przyjęty zestaw deskryptorów jest wzajemnie komplementarny.


Do analizy korelacji wybrano trzy nie-MFCC deskryptory: Centroid, F0 oraz PAR, uzupełnione o~trzy reprezentatywne współczynniki MFCC. Włączenie Centroidu, F0 i~PAR było kluczowe, ponieważ parametry te reprezentują trzy odrębne domeny wektora cech (odpowiednio: jasność widmową, wysokość dźwięku oraz dynamikę czasową); zestawienie ich w~jednej macierzy pozwala zweryfikować, czy rzeczywiście niosą wzajemnie niezależną informację, co stanowi bezpośredni test redundancji cech w~metryce opartej na odległości euklidesowej. Spośród współczynników MFCC wybrano MFCC\_1 (energia sygnału) oraz MFCC\_2 (globalna składowa barwy i~nachylenie widma), osiągające najwyższe wartości bezwzględne w~przestrzeni MFCC (odpowiednio $-303{,}938$ i~$98{,}238$), a~także MFCC\_9 jako przykład współczynnika wyższego rzędu, wybrany ze względu na najwyższą zmienność w~badanym zbiorze (odchylenie standardowe po normalizacji std~$= 0{,}286$). Prezentację ograniczono do tego podzbioru w~celu zachowania czytelności macierzy oraz uwypuklenia zależności między głównymi, interpretowalnymi parametrami akustycznymi. Ograniczenie to dotyczy jednak wyłącznie warstwy ilustracyjnej: w~analizie PCA składowe główne (PC1 i~PC2) wyznaczane są na podstawie pełnego wektora cech (Centroid, F0, PAR oraz MFCC\_1--MFCC\_13), dzięki czemu porównywanie próbek dźwiękowych uwzględnia łączną informację niesioną przez wszystkie deskryptory, a~nie jedynie wybrane współczynniki. Otrzymaną macierz korelacji przedstawiono na Rys.~\ref{fig:matrix}.



\vspace{1em}

\begin{figure}[H]
    \centering
    \includegraphics[width=0.9\textwidth]{images/correl-matrix.jpeg}
    \caption{Macierz korelacji}
    \label{fig:matrix}
\end{figure}

\begin{enumerate}[label=\roman*.]
    \item \textbf{Spectral Centroid, a~parametr F0}\\
    w~badanym zbiorze danych zależność między parametrem \textit{Spectral Centroid} a~częstotliwością podstawową F0 jest słaba (współczynnik korelacji wynosi $0{,}30$). Wynik ten wskazuje, że jasność dźwięku (percepcyjna barwa) nie jest bezpośrednio determinowana przez fizyczną wysokość tonu podstawowego. Analiza próbek wykazała obecność sygnałów o~unikalnej charakterystyce, takich jak \textit{tawny\_owl.mp3} oraz \textit{duck.mp3}, które mimo niskiej wysokości dźwięku (niskie wartości F0) charakteryzują się bardzo wysokim centroidem. Potwierdza to, że w~przypadku tych dźwięków energia spektralna jest przesunięta ku wyższym częstotliwościom ze względu na ich specyficzną, bogatą w~harmoniczne lub szumy strukturę, co czyni te parametry niezależnymi i~wartościowymi diagnostycznie źródłami informacji. Z~punktu widzenia percepcji dzwięku,  można mieć dźwięki niskie (niskie F0), które są jednocześnie "jasne" (wysoki centroid) przez dużą ilość składowych wysokoczęstotliwościowych. Stąd, tworząc metrykę podobieństwa, oba parametry są potrzebne.

    \item \textbf{Relacja Centroid, a~MFCC\_1 i~MFCC\_2} \\
    Najsilniejsze zależności zaobserwowano między jasnością dźwięku (centroid) a~składowymi barwy. Odnotowano istotną ujemną korelację z~MFCC\_1 ($-0.63$) oraz MFCC\_2 ($-0.62$). Sugeruje to, że wraz ze wzrostem jasności widmowej (centroidu), wartości MFCC\_1 i~MFCC\_2 maleją. w~praktyce oznacza to, że jaśniejsze dźwięki w~zbiorze mają specyficzny, powtarzalny profil energetyczny i~barwowy.

    % \item \textbf{Niezależność parametru F0} \\
    % w~zbiorze danych zależność między centroidem a~f0 ($0.30$) jest umiarkowanie słaba. Świadczy to o~tym, że jasność dźwięku (barwa) nie jest bezpośrednio determinowana przez wysokość tonu podstawowego. w~zbiorze znajdują się próbki, które mimo niskiej wysokości dźwięku (niskie F0) mogą być bardzo jasne (wysoki centroid), są to \textit{tow\_owl.mp3} oraz \textit{duck.mp3}.

    \item \textbf{Unikalność parametru MFCC\_9} \\
    Zaobserwowano silną ujemną korelację między MFCC\_9 a~energią całkowitą sygnału (MFCC\_1), wynoszącą $-0{,}66$. Oznacza to, że parametr ten ściśle reaguje na zmiany głośności i~natężenia dźwięku. Jednocześnie MFCC\_9 wykazuje niemal całkowitą niezależność od dynamiki wewnętrznej sygnału (parametr PAR, korelacja $0{,}10$). Brak związku z~impulsowym charakterem dźwięku sugeruje, że współczynnik ten opisuje cechy barwowo-teksturalne, a~nie chwilowe skoki amplitudy. Występuje umiarkowana korelacja MFCC\_9 z~częstotliwością podstawową (F0) na poziomie $0{,}33$. Wynik ten potwierdza, że informacja niesiona przez parametr MFCC\_9 jest w~dużym stopniu niezależna od tego, czy dźwięk jest wysoki, czy niski. Dzięki temu MFCC\_9 staje się narzędziem pozwalającym na obiektywną ocenę charakteru źródła dźwięku, bez wpływu jego wysokości.
    
    \item \textbf{Stabilność dynamiki (PAR)} \\
    Parametr PAR wykazuje najniższe korelacje z~większością pozostałych cech, szczególnie z~MFCC\_9 ($0.10$) i~MFCC\_2 ($-0.16$). Potwierdza to, że dynamika wewnętrzna próbek (stosunek szczytu do średniej) jest cechą niemal całkowicie niezależną od barwy dźwięku. Pozwala to na skuteczne rozróżnienie próbek pod względem ich "charakteru uderzeniowego" lub nasycenia, niezależnie od tego, czy dźwięk jest wysoki, czy niski. w~praktyce uzasadnia to wykorzystnie PAR w~wektorze cech jako osobnego wymiaru: pozwala rozróżniać próbki o~podobnej wysokości i~barwie, ale odmiennej "impulsowości" lub rozkładzie energii w~czasie.
\end{enumerate}

Rzutowanie próbek na płaszczyznę wyznaczoną przez dwie pierwsze główne składowe (PC1 i~PC2) przedstawiono na Rys.~\ref{fig:pc12}.


\begin{figure}[H]
    \centering
    \includegraphics[width=0.8\textwidth]{images/pca_normas.png}
    \caption{Cechy PC1 oraz PC2, elementy zielone z~czarną ramką oznaczają dzwięki pochodzące z~oryginalnego badania\cite{MadhavanWrobel2024MLT}.}
    \label{fig:pc12}
\end{figure}



\subsection{Część B: Weryfikacja metryki z~wynikami oceny klinicznej}


W~celu empirycznej walidacji opracowanej metryki podobieństwa obliczone odległości euklidesowe zestawiono z~wynikami badania percepcji słuchowej.

Dla czterech par dzwiękowych wyznaczono odległości euklidesowe w~16-wymiarowej znormalizowanej przestrzeni cech i~zestawiono je z~odsetkiem poprawnych identyfikacji liczby dźwięków w~badaniu Madhavan et al.~\cite{MadhavanWrobel2024MLT} (Tab.~\ref{tab:mlt_walidacja}).

\begin{table}[H]
\centering
\begin{tabular}{llcc}
\toprule
Para bodźców & Odległość euklidesowa & Rozróżnialność \% prawidłowych odpowiedzi \\
\midrule
cow.mp3 - horse.mp3     & 0{,}714 & 20{,}5\% \\
cow.mp3 - lion.mp3    & 0{,}942 & 92{,}9\% \\
donkey.mp3 - horse.mp3   & 1{,}093 & 55{,}6\% \\
sheep.mp3 - frog.mp3  & 1{,}736 & 44{,}4\% \\
\bottomrule
\end{tabular}
\caption{Zestawienie odległości euklidesowych między parami bodźców dźwiękowych z~odsetkiem poprawnych odpowiedzi w~próbach dwuelementowych badania MLT~\cite{MadhavanWrobel2024MLT}.}
\label{tab:mlt_walidacja}
\end{table}

Dla trzech spośród czterech analizowanych par zaobserwowano zgodność z~przewidywaniami modelu: mniejszy dystans w~przestrzeni cech koreluje z~niższą dokładnością identyfikacji liczby jednocześnie prezentowanych dźwięków. Para \textit{cow.mp3-horse.mp3}, charakteryzująca się najniższą odległością euklidesową ($d = 0{,}714$), okazała się najtrudniejsza do rozróżnienia dla słuchaczy ($20{,}5\%$ poprawnych odpowiedzi). Para \textit{cow.mp3-lion.mp3}, o~wyraźnie większym dystansie ($d = 0{,}942$), była identyfikowana poprawnie w~niemal $93\%$ przypadków. Pośrednia para \textit{donkey.mp3-horse.mp3} ($d = 1{,}093$) uzyskała pośrednią dokładność $55{,}6\%$, co w~pełni wpisuje się w~przewidziany gradient trudności.

Wyjątek stanowi para \textit{sheep.mp3-frog.mp3} ($d = 1{,}736$), która mimo największego dystansu akustycznego osiągnęła jedynie $44{,}4\%$ poprawnych identyfikacji. 



\section{Dyskusja}

Pytanie, jak porównać dwa dźwięki, jest pozornie proste. W~praktyce podobieństwo akustyczne i~podobieństwo percepcyjne są konstruktami wielowymiarowymi i~nie mogą być wiarygodnie opisane przez pojedynczy parametr sygnału. Już klasyczne badania nad barwą dźwięku wykazały, że ocena podobieństwa opiera się równocześnie na kilku częściowo niezależnych cechach spektralnych i~czasowych. Grey wykorzystywał wielowymiarowe skalowanie ocen podobieństwa brzmień instrumentów, natomiast Lakatos wykazał, że do głównych wymiarów percepcyjnych należą m.in. centroid widmowy i~czas narastania sygnału~\cite{Grey1977Timbre,Lakatos2000Timbre}. Późniejsze badania potwierdziły, że percepcyjna przestrzeń barwy może być opisana kombinacją kilku cech akustycznych, a~nie pojedynczą zmienną~\cite{Elliott2013TimbreDimensions}.

W~odniesieniu do dźwięków środowiskowych szczególnie istotny jest wynik pracy Gygi i~wsp., w~której analizowano 50 naturalnych dźwięków. Autorzy wykazali, że oceny podobieństwa tworzą wielowymiarową strukturę, a~położenie bodźców w~tej przestrzeni można częściowo przewidywać na podstawie kombinacji takich cech, jak harmoniczność, udział ciszy czy głębokość modulacji~\cite{Gygi2007EnvironmentalSounds}. Wynik ten ma bezpośrednie znaczenie dla niniejszej pracy, ponieważ analizowany materiał dźwiękowy również obejmuje złożone, semantycznie rozpoznawalne odgłosy, a~nie laboratoryjne tony czyste. Oznacza to, że podobieństwo pomiędzy takimi bodźcami powinno być traktowane jako wynik współdziałania wielu właściwości sygnału, a~nie jako funkcja samej wysokości tonu czy pojedynczej miary widmowej.

Współczesne metody analizy dźwięku umożliwiają odwrócenie klasycznego paradygmatu badań psychofizycznych: zamiast najpierw pytać słuchacza, które dźwięki są podobne, można najpierw zbudować numeryczną reprezentację każdego bodźca, a~następnie obliczyć odległość pomiędzy reprezentacjami. Rozwój zestawów deskryptorów, takich jak \textit{Timbre Toolbox} (zestaw narzędzi do automatycznego opisu barwy i~innych właściwości dźwięku za pomocą wielu parametrów akustycznych), pozwolił ujmować strukturę sygnału za pomocą wielu parametrów spektralnych, energetycznych i~czasowych~\cite{Peeters2011TimbreToolbox}. Współczynniki MFCC stały się natomiast standardowym sposobem zwartego opisu kształtu obwiedni widmowej, szeroko stosowanym w~analizie mowy, muzyki, identyfikacji źródeł oraz klasyfikacji dźwięków. Należy podkreślić, że ani zastosowanie MFCC, ani wykorzystanie odległości euklidesowej nie stanowią same w~sobie nowości metodologicznej. Już Helén i~Virtanen wykorzystywali miary odległości pomiędzy reprezentacjami cech w~systemie \textit{audio query-by-example}, porównując m.in. odległość euklidesową z~miarami opartymi na rozkładach prawdopodobieństwa~\cite{HelenVirtanen2010AudioQuery}. De Camargo i~wsp. zastosowali MFCC oraz miary odległości w~automatycznej identyfikacji odgłosów zwierząt~\cite{DeCamargo2017ProtaxSound}, z~kolei Siedenburg i~Müllensiefen modelowali jawne oceny podobieństwa krótkich fragmentów muzycznych za pomocą zestawów deskryptorów barwy i~MFCC, wykorzystując odległości pomiędzy cechami jako predyktory ocen słuchaczy~\cite{Siedenburg2017TimbreSimilarity}. Badania te potwierdzają, że wektorowa reprezentacja sygnału jest uzasadnionym sposobem ilościowego ujmowania podobieństwa, ale jednocześnie wyznaczają granicę, poza którą należy poszukiwać właściwej innowacyjności obecnego podejścia.

Odrębność niniejszej pracy nie polega zatem na zaproponowaniu nowej klasy deskryptorów ani nowej matematycznej definicji dystansu. Polega na przypisaniu tej odległości innej roli eksperymentalnej. W~typowych systemach rozpoznawania lub wyszukiwania dźwięków dystans pomiędzy wektorami służy jako kryterium klasyfikacji: odpowiedź ma wskazać, jaki to dźwięk albo który element bazy jest najbardziej podobny do dźwięku, o~który pytamy. W~badaniach psychofizycznych odległość w~przestrzeni cech bywa z~kolei odnoszona do jawnej oceny podobieństwa, wydawanej przez słuchacza. W~obecnym badaniu odległość akustyczna jest natomiast traktowana jako niezależna cecha konstrukcji bodźców i~potencjalny czynnik wpływający na wykonanie zadania percepcyjno-poznawczego. Jest to istotne, ponieważ uczestnicy pierwotnego badania, opisanego przez Madhavan i~wsp.~\cite{MadhavanWrobel2024MLT}, nie odpowiadali na pytanie, czy dwa dźwięki są do siebie podobne. Ich zadaniem było rozpoznanie liczby jednocześnie prezentowanych obiektów słuchowych. Podobieństwo akustyczne nie było więc bezpośrednim przedmiotem odpowiedzi, lecz mogło działać jako ukryty czynnik utrudniający segregację sceny słuchowej. W~tym ujęciu pytanie „jak porównać dwa dźwięki?” ma znaczenie metodologiczne: chodzi nie tylko o~opis ich podobieństwa, ale o~ocenę, czy para bodźców jest na tyle bliska akustycznie, że może zmieniać trudność testu niezależnie od procesów poznawczych, które test ma mierzyć.

Takie zastosowanie metryki ma praktyczne znaczenie dla projektowania materiału w~badaniach audiologicznych i~kognitywnych. Dwa bodźce mogą należeć do odmiennych kategorii semantycznych, a~mimo to pozostawać bliskie w~wielowymiarowej przestrzeni cech. Jeżeli są prezentowane jednocześnie, ich rozdzielenie może być utrudnione przez podobieństwo struktury widmowej, wysokości, dynamiki lub innych właściwości sygnału. W~konsekwencji wynik testu może odzwierciedlać zarówno badaną zdolność poznawczą, jak i~właściwości samego materiału dźwiękowego. Obiektywna analiza odległości między bodźcami może więc pełnić funkcję kontroli jakości zestawu testowego i~pomagać w~identyfikacji par nadmiernie podobnych jeszcze przed rozpoczęciem badania klinicznego.

Uzyskane wyniki dostarczają wstępnego wsparcia dla takiej interpretacji. W~trzech spośród czterech par, dla których dostępne były dane z~badania Madhavan i~wsp.~\cite{MadhavanWrobel2024MLT}, obserwowano kierunek zgodny z~założeniem, że mniejszej odległości w~przestrzeni cech może towarzyszyć większa trudność zadania. Para \textit{cow.mp3-horse.mp3}, charakteryzująca się najmniejszym dystansem spośród tych czterech porównań ($d = 0{,}714$), była jednocześnie parą o~najniższym odsetku prawidłowych odpowiedzi ($20{,}5\%$). Para \textit{cow.mp3-lion.mp3}, bardziej oddalona w~przestrzeni cech ($d = 0{,}942$), wiązała się z~$92{,}9\%$ poprawnych odpowiedzi, natomiast \textit{donkey.mp3-horse.mp3} ($d = 1{,}093$) zajmowała pozycję pośrednią pod względem uzyskanych wyników. Taki układ jest zgodny z~hipotezą, że akustyczna bliskość bodźców może zwiększać obciążenie procesów segregacji słuchowej. Interpretacja ta musi jednak pozostać ostrożna. Para \textit{sheep.mp3-frog.mp3} stanowiła wyjątek: mimo największej odległości euklidesowej ($d = 1{,}736$) uzyskano jedynie $44{,}4\%$ odpowiedzi prawidłowych. Wynik ten pokazuje, że pojedyncza globalna miara odległości nie może być traktowana jako kompletny predyktor trudności percepcyjnej. Na wykonanie zadania mogą wpływać również właściwości czasowe sygnału, synchronia początku bodźców, modulacja amplitudy i~częstotliwości, różnice w~czasie trwania, cechy semantyczne i~rozpoznawalność źródła, a~także mechanizmy uwagi i~organizacji sceny słuchowej. Z~tego względu obecne dane należy traktować jako pilotażową walidację koncepcji, a~nie jako ustalenie uniwersalnego progu podobieństwa.

Warto również skorygować interpretację rozszerzonego zestawu 24 bodźców. Minimalna odległość w~pierwotnym zbiorze MLT wynosiła $0{,}714$ dla pary \textit{cow.mp3-horse.mp3}, natomiast w~zbiorze rozszerzonym znaleziono jeszcze bliższą parę \textit{owl.mp3-tawny\_owl.mp3} ($d = 0{,}502$). Oznacza to, że rozszerzenie zbioru rzeczywiście wprowadziło parę o~większym podobieństwie akustycznym, a~nie -- jak można byłoby błędnie wnioskować -- wyłącznie poszerzyło pokrycie przestrzeni cech bez zwiększenia redundancji. Z~punktu widzenia celu pracy, taki wynik pokazuje, że zaproponowana procedura może wykrywać potencjalnie problematyczne pary bodźców, zanim zostaną one zastosowane w~badaniu percepcyjnym.

\subsection{Znaczenie doboru cech i~ograniczenia interpretacyjne}

Dobór centroidu widmowego, F0, PAR oraz MFCC tworzy relatywnie prosty i~interpretowalny wektor obejmujący różne aspekty sygnału. Zaletą takiego rozwiązania jest możliwość sprowadzenia złożonego opisu akustycznego do jednej miary odległości przy zachowaniu informacji o~wielu cechach jednocześnie. Wyniki analizy korelacji wskazują, że część parametrów wnosi względnie odrębną informację -- przykładowo zależność pomiędzy F0 i~centroidem widmowym była słaba. Jest to argument za łączeniem deskryptorów, zamiast redukowania podobieństwa do pojedynczego wymiaru. Należy jednak unikać nadmiernie literalnej interpretacji poszczególnych współczynników MFCC. Współczynniki te opisują kolejne składowe obwiedni widmowej w~skali melowej, przy czym wraz ze wzrostem numeru współczynnika reprezentowane są coraz drobniejsze szczegóły kształtu widma~\cite{Siedenburg2017TimbreSimilarity}. Nie istnieje uniwersalne przyporządkowanie typu „MFCC\_5 = określony rezonans” czy „MFCC\_9 = mikrostruktura” niezależne od materiału i~parametrów ekstrakcji. Z~tego samego powodu wysoka średnia wartość znormalizowanego MFCC nie jest równoznaczna z~wysoką zdolnością dyskryminacyjną tej cechy. Ocena rzeczywistego znaczenia poszczególnych wymiarów wymagałaby np. analizy ważności cech, walidacji krzyżowej modeli z~usuwaniem kolejnych parametrów albo zastosowania metod redukcji wymiarowości i~regresji na większym zbiorze danych.

Drugim ograniczeniem jest sposób agregowania informacji w~czasie. W~obecnej procedurze dla całego nagrania uzyskuje się uśrednione wartości cech. Takie podejście jest proste i~stabilne, lecz traci informacje o~dynamice sygnału: początku dźwięku, obwiedni, modulacji, zmianach widma oraz kolejności zdarzeń. Lakatos wykazał znaczenie czasu narastania dla percepcji barwy~\cite{Lakatos2000Timbre}, a~praca Siedenburga i~Müllensiefena pokazuje, że cechy spektroczasowe i~pochodne MFCC mogą wnosić istotną informację przy modelowaniu podobieństwa~\cite{Siedenburg2017TimbreSimilarity}. W~przyszłości zasadne byłoby więc porównanie obecnego modelu z~reprezentacją uwzględniającą delta-MFCC, parametry obwiedni czasowej lub metody dopasowania sekwencji, np. \textit{dynamic time warping}.

Odległość euklidesowa ma w~tym projekcie ważną zaletę: jest prosta, jednoznaczna i~łatwa do interpretacji. Po normalizacji cech każdy parametr uczestniczy w~obliczeniu dystansu w~porównywalnej skali, co pozwala uzyskać przejrzystą miarę globalną. Jednocześnie wybór ten implikuje założenie o~równym znaczeniu wszystkich wymiarów i~o~geometrycznie izotropowej przestrzeni cech. Założenie to nie musi być zgodne z~rzeczywistą wagą percepcyjną poszczególnych parametrów.

Literatura dotycząca wyszukiwania i~modelowania podobieństwa audio wskazuje, że skuteczność konkretnej metryki zależy od rodzaju materiału i~celu analizy~\cite{HelenVirtanen2010AudioQuery,Siedenburg2017TimbreSimilarity}. Kolejnym etapem badań powinno więc być porównanie dystansu euklidesowego z~innymi miarami, np. odległością kosinusową, odległością Mahalanobisa lub metrykami uczonymi na danych behawioralnych. Szczególnie interesujące byłoby wyznaczenie wag cech na podstawie większego zbioru odpowiedzi MLT, zamiast przyjmowania jednakowego udziału każdego wymiaru. Obecny model należy zatem traktować jako celowo prostą i~interpretowalną metrykę bazową, stanowiącą punkt odniesienia dla bardziej złożonych rozwiązań.

Najbardziej praktycznym zastosowaniem przedstawionej procedury jest możliwość wstępnej standaryzacji materiału dźwiękowego. W~testach, których celem jest ocena procesów poznawczych, pamięciowych lub integracji percepcyjnej, bodźce powinny być dobrane w~taki sposób, aby niezamierzona bliskość akustyczna nie stała się dominującym źródłem trudności. Zastosowanie macierzy odległości dla całego zestawu pozwala wskazać pary szczególnie bliskie, kontrolować rozkład trudności oraz budować zestawy alternatywne o~podobnym profilu akustycznym. Jest to potencjalnie użyteczne nie tylko dla MLT, ale także dla innych procedur audiologicznych i~eksperymentów nad percepcją słuchową. Jednocześnie obecne badanie ma charakter eksploracyjny. Liczba bodźców jest niewielka, a~walidacja behawioralna obejmuje jedynie cztery pary. Na tej podstawie nie można wyznaczyć wiarygodnych progów klasyfikujących dźwięki jako „podobne”, „pośrednie” lub „różne”. W~kolejnych badaniach konieczna będzie walidacja większej liczby par, analiza prób obejmujących więcej niż dwa jednoczesne bodźce oraz bezpośrednie porównanie miary akustycznej zarówno z~jawnymi ocenami podobieństwa, jak i~wynikami zadania MLT. Takie rozwiązanie pozwoliłoby rozdzielić trzy poziomy: podobieństwo fizyczne sygnału, świadomie raportowane podobieństwo percepcyjne oraz rzeczywisty wpływ podobieństwa na wykonanie zadania. W~tym kontekście proponowana metoda nie powinna być przedstawiana jako uniwersalny sposób rozstrzygania, czy dwa dźwięki są „takie same” lub „różne”. Jej wartość polega raczej na dostarczeniu obiektywnego, powtarzalnego i~łatwego do audytu wskaźnika bliskości akustycznej, który można wykorzystać jako dodatkową zmienną przy projektowaniu i~interpretacji testów słuchowych. Takie ujęcie łączy analizę sygnału z~psychofizyką i~zastosowaniem klinicznym: od pytania, jakie cechy zawiera dźwięk, przechodzi do pytania, w~jakim stopniu właściwości bodźca mogą współkształtować zachowanie słuchacza.

Podsumowując, wcześniejsze badania jednoznacznie pokazują, że ilościowe porównywanie dźwięków za pomocą wielowymiarowych reprezentacji i~miar odległości ma ugruntowane podstawy metodologiczne. Nowość proponowanego podejścia należy więc lokować nie w~samych narzędziach analizy sygnału, lecz w~ich zastosowaniu do kontroli materiału bodźcowego i~interpretacji wykonania zadania percepcyjnego. Wstępna zgodność pomiędzy małą odległością akustyczną a~niższą poprawnością w~części par MLT sugeruje, że wielocechowa metryka może pomagać identyfikować bodźce niosące dodatkowe ryzyko interferencji. Jednocześnie obecność wyniku odstającego oraz mała liczba dostępnych porównań wskazują, że zależność ta wymaga dalszej, prospektywnej walidacji. W~takim ujęciu przedstawiona procedura stanowi nie końcowy model percepcji dźwięku, lecz praktyczne i~przejrzyste narzędzie do zadania fundamentalnego pytania: jak porównać dwa dźwięki w~sposób istotny dla zachowania słuchacza?

\section{Wnioski}

Na podstawie przeprowadzonych analiz sformułowano następujące wnioski:

\begin{enumerate}[label=\roman*.]
    \item Opracowano 16-wymiarową metrykę podobieństwa akustycznego opartą na znormalizowanych wektorach cech (centroid widmowy, F0, PAR oraz MFCC\_1-MFCC\_13), umożliwiającą obiektywne, powtarzalne porównywanie próbek dźwiękowych niezależnie od pierwotnych jednostek miary poszczególnych parametrów.

    \item Analiza korelacji wykazała niską zależność między centroidem widmowym a~częstotliwością podstawową ($r = 0{,}30$), co potwierdza komplementarność wybranych deskryptorów: jasność barwy i~wysokość tonu niosą niezależną informację i~łącznie zwiększają zdolność różnicowania podobieństwa akustycznego.

    \item Spośród wszystkich współczynników MFCC największą zdolność różnicującą wykazał MFCC\_5 (najwyższa średnia znormalizowana, $\overline{x} = 0{,}617$), natomiast parametr PAR okazał się niemal całkowicie niezależny od cech barwowych ($r \leq 0{,}16$ względem wszystkich współczynników MFCC), co uzasadnia jego włączenie do wektora jako osobnego wymiaru opisu dynamiki sygnału.

    \item Opracowana procedura dostarcza praktycznego narzędzia do projektowania zestawów bodźców akustycznych o~z~góry zdefiniowanym i~kontrolowanym stopniu trudności percepcyjnej, co stanowi bezpośrednie zastosowanie w~standaryzacji materiału dźwiękowego wykorzystywanego w~testach audiologicznych i~kognitywnych.
\end{enumerate}

\clearpage
\bibliography{bibliography}


\clearpage
\appendix
\section{Wektory cech akustycznych}


Wektor cech akustycznych przed normalizacją min-max, wykorzystanych w~artykule.

\begin{tabular}{|c|c|c|c|}\hline
Name & F0 & PAR & Centroid \\\multicolumn{4}{|c|}{MFCC } \\ \hline
chaffinch & 1397.990 & 21.744 & 4473.348 \\\multicolumn{4}{|c|}{-379.172 84.621 -16.691 34.242 -16.974 -18.007 -17.783 18.617 -12.605 -0.168 -3.704 -1.220 -2.195 } \\ \hline
chicken & 930.343 & 17.700 & 2926.691 \\\multicolumn{4}{|c|}{-327.237 82.999 -38.592 -33.489 -32.739 -16.701 -22.290 -11.800 4.410 12.632 5.750 2.132 0.797 } \\ \hline
cow & 875.505 & 18.591 & 1654.120 \\\multicolumn{4}{|c|}{-394.512 137.184 -21.177 -26.304 -1.867 9.404 0.484 -2.328 2.940 1.853 -4.083 -4.867 0.937 } \\ \hline
dog & 501.156 & 10.051 & 2209.468 \\\multicolumn{4}{|c|}{-210.405 140.009 -0.183 11.668 -14.332 0.444 -15.602 -17.633 -22.330 -8.135 -1.243 -5.612 1.156 } \\ \hline
donkey & 1254.720 & 16.949 & 3438.254 \\\multicolumn{4}{|c|}{-210.732 119.803 -76.757 1.164 -57.688 24.694 -14.443 14.486 -12.009 0.844 12.634 -10.031 1.196 } \\ \hline
duck & 809.362 & 27.325 & 4146.062 \\\multicolumn{4}{|c|}{-214.189 113.623 -54.584 34.472 -32.769 29.992 -17.735 16.085 -13.393 13.294 -6.171 8.963 -3.980 } \\ \hline
elephant & 922.155 & 13.935 & 1530.462 \\\multicolumn{4}{|c|}{-260.300 162.440 -58.624 -32.465 16.992 -23.840 -13.287 14.540 -10.512 -11.306 6.717 -6.099 -4.969 } \\ \hline
frog & 1580.831 & 18.955 & 5692.735 \\\multicolumn{4}{|c|}{-483.222 33.491 -25.037 -1.014 1.969 0.867 0.394 1.640 6.719 6.226 -1.556 -4.698 -0.447 } \\ \hline
frog2 & 1016.737 & 23.145 & 3016.866 \\\multicolumn{4}{|c|}{-392.174 131.438 -36.168 2.720 9.223 13.644 0.556 7.303 7.392 10.260 17.142 1.723 1.871 } \\ \hline
gorilla & 786.202 & 17.734 & 1330.571 \\\multicolumn{4}{|c|}{-82.167 96.097 -65.205 12.497 -24.149 6.374 -19.602 -0.698 -23.692 -3.660 -22.783 13.400 -6.624 } \\ \hline
great\_tit & 1421.528 & 18.490 & 3894.647 \\\multicolumn{4}{|c|}{-401.115 71.831 -31.917 44.812 5.723 -18.835 8.275 9.700 -3.537 -6.278 5.170 4.607 -4.099 } \\ \hline
hawk & 2420.861 & 15.379 & 2490.972 \\\multicolumn{4}{|c|}{-251.887 103.621 -128.249 12.337 -16.913 0.574 -9.382 -23.244 18.037 1.230 13.741 -0.940 8.283 } \\ \hline
horse & 1067.942 & 17.368 & 2011.143 \\\multicolumn{4}{|c|}{-297.051 119.692 -53.103 -3.729 -8.849 7.642 -5.801 -3.085 -13.738 -5.808 -6.445 -3.819 -1.159 } \\ \hline
lion & 437.437 & 12.647 & 1858.199 \\\multicolumn{4}{|c|}{-294.076 104.998 -4.434 26.838 -4.114 10.453 -3.604 8.942 -4.704 4.246 0.125 -1.112 -0.304 } \\ \hline
monkey & 1710.385 & 22.929 & 3530.151 \\\multicolumn{4}{|c|}{-394.827 53.336 -97.621 -15.711 -11.431 7.840 -5.534 18.661 11.501 3.127 -7.646 -2.045 -11.244 } \\ \hline
owl & 622.496 & 13.560 & 4539.256 \\\multicolumn{4}{|c|}{-429.998 47.405 10.065 -2.331 -8.078 -7.900 -2.792 10.686 15.895 15.074 8.192 2.367 -1.195 } \\ \hline
sheep & 1279.340 & 19.892 & 3325.403 \\\multicolumn{4}{|c|}{-189.880 109.772 -88.828 6.104 -25.419 15.384 -30.839 -3.600 -21.130 9.763 -1.861 0.301 12.754 } \\ \hline
tawny\_owl & 456.731 & 17.654 & 5257.504 \\\multicolumn{4}{|c|}{-447.315 38.390 15.328 5.323 1.278 -0.280 0.894 6.429 9.165 9.383 6.560 4.190 1.832 } \\ \hline
tundra\_swan & 784.416 & 18.319 & 2701.883 \\\multicolumn{4}{|c|}{-230.075 125.439 -51.170 5.210 -21.519 -6.252 -26.662 -14.046 -10.482 5.757 13.905 9.226 7.296 } \\ \hline
wild\_boar & 1085.865 & 22.563 & 3304.438 \\\multicolumn{4}{|c|}{-268.061 126.635 -52.249 20.793 3.746 11.124 -4.933 -14.149 7.141 -1.544 -3.149 2.109 -6.019 } \\ \hline
wolf1 & 510.134 & 15.807 & 2421.858 \\\multicolumn{4}{|c|}{-231.524 144.509 -9.122 21.388 -25.433 -2.456 -14.032 -4.206 -9.847 -1.886 -8.958 2.777 -3.800 } \\ \hline
wolf2 & 551.040 & 16.948 & 2131.952 \\\multicolumn{4}{|c|}{-334.230 131.805 41.666 -0.186 -5.520 5.895 1.657 5.204 -0.917 2.270 -2.057 7.741 11.357 } \\ \hline
blackbird & 2240.552 & 20.041 & 4619.491 \\\multicolumn{4}{|c|}{-381.544 39.774 -62.417 10.515 18.136 4.618 -10.540 11.598 1.926 7.897 1.459 2.245 4.149 } \\ \hline
cat & 1100.218 & 17.867 & 1513.415 \\\multicolumn{4}{|c|}{-188.818 38.797 -19.753 17.932 -11.072 15.299 1.157 43.015 0.418 4.087 -22.761 9.654 4.120 } \\ \hline
\end{tabular}





\end{document}
